\section{Limitations and Future Directions}
\label{sec:limitations}

\subsection{Current Limitations}
\label{ssec:limitations_current}

\paragraph{Annotation subjectivity.} Causal reasoning inherently admits multiple
valid interpretations. Although our three-stage quality control pipeline and
inter-annotator agreement metrics ($\kappa = 0.81$--$0.88$) indicate strong consensus,
some ambiguous scenarios may have been resolved with a systematic bias reflecting
the training received by annotators at a specific annotation centre. Future work
should investigate cross-cultural variation in causal attribution as an independent
research problem.

\paragraph{Geographic imbalance.} While \lucid{} spans six cities, Shanghai
contributes 33.6\% of all scenarios --- a consequence of the largest fleet deployment
in that city during the initial collection phase. Future extensions should target
additional cities in Africa, Latin America, and South/Southeast Asia, which are
significantly underrepresented in existing AD datasets.

\paragraph{LiDAR-camera synchronisation.} The rooftop LiDAR operates at 20~Hz while
cameras operate at 10~FPS; sub-frame temporal misalignment of up to 50~ms may
introduce minor artefacts in LiDAR-projected depth images used for annotation.
A next-generation collection platform with hardware-triggered frame synchronisation
is under development.

\paragraph{Simulation gap.} All \lucid{} scenarios are collected from real-world
driving; the dataset does not include synthetic data. While this ensures ecological
validity, it limits the coverage of truly extreme scenarios (\eg{} vehicle rollovers,
bridge collapse) that are too dangerous to collect in the field. A synthetic
augmentation study using CARLA-simulated edge cases is planned.

\paragraph{Language coverage.} All annotations are authored in English. Extending
\lucid{} to Mandarin Chinese, German, and Singaporean English would support
evaluation of multilingual LLMs and mitigate potential bias toward English-centric
models.

\subsection{Future Directions}
\label{ssec:limitations_future}

\begin{itemize}[leftmargin=2em, noitemsep]
  \item \textbf{Causal structure prediction}: Training models to generate CEG graphs
        from raw sensor data, rather than text-only QA, would enable richer causal
        representation learning.
  \item \textbf{Active learning for annotation}: The annotation cost of \lucid{}
        ($\sim$45,000 person-hours) motivates research into active learning strategies
        that prioritise high-value scenarios for human annotation.
  \item \textbf{Embodied agent evaluation}: Future work will extend the IFN task to
        closed-loop evaluation in CARLA, enabling end-to-end assessment of instruction-conditioned
        driving policies.
  \item \textbf{Temporal counterfactuals}: Current counterfactual pairs consider
        single-step interventions; structured temporal counterfactuals (\eg{} ``What if
        the light had changed 5 seconds earlier?'') would test multi-step causal reasoning.
\end{itemize}
